\documentclass[nonacm, acmsmall]{acmart}
\usepackage{note}

\title{Research note template}

\author{Your Name}
\affiliation{%
	\institution{Carnegie Mellon University}
	\city{Pittsburgh}
	\state{Pennsylvania}
	\country{USA}
}

\begin{document}

\maketitle

\tableofcontents

\section{Introduction}
\label{sec:introduction}

Use this document as a template for research notes. Representative citation of a cool paper \citep{Miller:VOS:2024}.

\section{Notation and results}
\label{sec:results}

Notation macros belong in \texttt{note.sty}, so that changing a symbol is one edit. Equations are typeset with the packages and macros \texttt{mathdefs} provides,
%
\begin{equation}
	\label{equ:example}
	\E{f \paren{\bm{x}}} = \int_{\R^3} f \paren{\bm{x}} \, p \paren{\bm{x}} \ud \bm{x},
\end{equation}
%
and \texttt{boxdefs} provides environments for the statements you want to stand out.

\begin{prp}{An example proposition}{example}
	State the claim here. The environment numbers propositions, definitions, and lemmas in a single sequence, and labels this one \texttt{pro:example}.
\end{prp}

\Cref{pro:example} follows from \cref{equ:example}, which is the sort of cross-reference \texttt{cleveref} writes out for you.

\bibliographystyle{ACM-Reference-Format}
\bibliography{note}

\end{document}
